\section{Convergence against a non-pluripolar measure}
\label{sec:psh-fixed-measure}

A common bound on the potentials allows $L^1$ convergence for smooth volume to be transferred to a fixed finite measure that does not charge pluripolar sets. Bounded exponentials extend this statement to upper-bounded families.

Cegrell's fixed-measure convergence lemma says that a uniformly bounded sequence of plurisubharmonic functions converging in $L^1_{\mathrm{loc}}$ converges in $L^1(\nu)$ on compact subsets for each finite non-pluripolar measure $\nu$ there. The usual formulation allowing extraction gives the same conclusion: apply it to every subsequence, extract almost-everywhere convergence for smooth volume, and use the subsequence criterion for convergence. This is the classical fixed-measure argument used in \cite[Section 11.1.2]{GZ17}. We retain this local lemma as a standard input.

\Needspace{15\baselineskip}
\begin{lemma}\label{lem:fixed-npp-integrals}
Let $\nu$ be a finite non-pluripolar measure on $X$. If uniformly bounded quasi-plurisubharmonic functions $a_j,a$ have a common smooth lower curvature bound and $a_j\to a$ in $L^1(X)$, then
\[
 \int_X|a_j-a|\,\mathrm{d}\nu\longrightarrow0.
\]
If instead $u_j,u\in\PSH(X,\theta)$ have a common upper bound and $u_j\to u$ in $L^1(X)$, then, for every $b>0$,
\[
 \int_X|\mathrm{e}^{bu_j}-\mathrm{e}^{bu}|\,\mathrm{d}\nu\longrightarrow0.
\]
\end{lemma}
\begin{proof}
For the first statement add a fixed smooth strictly plurisubharmonic correction on each of finitely many coordinate balls. The resulting functions are uniformly bounded and plurisubharmonic. Apply the local input on smaller compact balls covering $X$ and sum the integrals; the corrections cancel in the difference.

For the second statement choose $A\omega\geq\theta$. If $u_j\leq C$, the convex chain rule gives
$\ddc \mathrm{e}^{bu_j}\geq-bA\mathrm{e}^{bC}\omega$, also for unbounded $u_j$ by bounded truncation. Thus $\mathrm{e}^{bu_j}$ are uniformly bounded quasi-plurisubharmonic functions in one lower curvature class. Their $L^1(\mathrm{d}V)$ convergence follows from the Lipschitz bound $b\mathrm{e}^{bC}$ for the exponential on $(-\infty,C]$. Apply the first statement. We assign the value zero to the exponential at a pole.
\end{proof}

\begin{proposition}\label{prop:singular-no-uniform-limit}
Let $\theta$ be smooth and let $\varphi\in\PSH(X,\theta)$. Suppose $\varphi$ is unbounded below and $u_\lambda\sim\varphi$ for every $\lambda>0$. Then $\lambda u_\lambda$ cannot converge uniformly on $X$ to a finite constant as $\lambda\downarrow0$. Its distance to any finite constant is also infinite in $L^\infty(\mathrm{d}V)$ for each fixed $\lambda>0$.
\end{proposition}
\begin{proof}
Bounded equivalence preserves unboundedness below. Every nonempty strict sublevel set of the upper semicontinuous function $u_\lambda$ is open and has positive smooth volume. Thus $u_\lambda$ is essentially unbounded below, as is $\lambda u_\lambda$. The stated conclusions follow.
\end{proof}

A vanishing exponent gives convergence to one even when the normalized potentials vary.

\Needspace{17\baselineskip}
\begin{lemma}\label{lem:normalized-exponential-lp}\label{lem:normalized-exponential-npp}
Let $v_j\in\PSH(X,\theta)$ satisfy $\sup_Xv_j=0$, and let $\lambda_j>0$ decrease to zero. For every finite non-pluripolar measure $\mu$,
\[
 \int_X\left|\mathrm{e}^{\lambda_jv_j}-1\right|\,\mathrm{d}\mu\longrightarrow0.
\]
If $f\in L^p(\mathrm{d}V)$ for some $p>1$, then more precisely
\[
 \int_X\left|\mathrm{e}^{\lambda_jv_j}-1\right|\,|f|\,\mathrm{d}V
 \leq C_p\lambda_j\|f\|_{L^p(\mathrm{d}V)}.
\]
The constant depends on the fixed background, $p$ and $\mathrm{d}V$.
\end{lemma}
\begin{proof}
Choose a Hermitian form $\Omega\geq\theta$. By \cref{lem:global-capacity-estimates}, for every $T>1$,
\[
 \operatorname{Cap}_\Omega(\{v_j<-T\})\leq C/T
\]
uniformly in $j$. Since $v_j\leq0$, splitting at this sublevel gives
\[
 \int_X(1-\mathrm{e}^{\lambda_jv_j})\,\mathrm{d}\mu
 \leq\lambda_jT\mu(X)+\mu(\{v_j<-T\}).
\]
\Cref{lem:npp-uniform-capacity} makes the second term uniformly small when $T$ is large. For this fixed $T$, the first term tends to zero as $j\to\infty$. This proves the general assertion. The value of the exponential is zero at a pole; pole sets have zero $\mu$-measure.

For the quantitative assertion, \cref{lem:uniform-exponential}, applied in background $\Omega$, gives a uniform $L^q$ bound for $-v_j$, where $q=p/(p-1)$. H\"older's inequality and
\[
 0\leq1-\mathrm{e}^{\lambda_jv_j}\leq\lambda_j|v_j|
\]
give the claimed estimate for every $j$.
\end{proof}

A moving bounded-product measure can also be used if its potentials converge in capacity. The proof estimates the torsion terms in integration by parts.

\begin{lemma}\label{lem:moving-bounded-integrals}
Let $a_j,a,w_j,w\in\PSH(X,\omega)$ be uniformly bounded. If $a_j\to a$ in $L^1(X)$ and $w_j\to w$ in capacity, then, for $0\leq k\leq n$,
\[
 \int_X|a_j-a|\,\omega_{w_j}^{k}\wedge\omega^{n-k}
 \longrightarrow0.
\]
The same conclusion holds when the $a_j$ have any fixed smooth lower curvature bound.
\end{lemma}
\begin{proof}
By the subsequence criterion, we may first extract $a_j\to a$ almost everywhere for smooth volume. The upper tail envelopes
$A_j=(\sup_{\ell\geq j}a_\ell)^*$ decrease to $a$: the local Hartogs lemma gives the upper inequality for the limit, and almost-everywhere convergence gives the reverse one. Bounded monotone capacity convergence therefore implies
\[
 \alpha_j:=\max(a_j,a)\longrightarrow a\quad\text{in capacity}.
\]
The part $\alpha_j-a$ of $|a_j-a|=(\alpha_j-a)+(\alpha_j-a_j)$ tends to zero against the moving products by \cref{lem:bounded-mixed-capacity}.

Put $h_j=\alpha_j-a_j\geq0$. We compare its integrals for two indices $j,\ell$. The case $k=0$ is the assumed $L^1$ convergence. For $k\geq1$, put $z=w_j-w_\ell$ and
\[
 T=\sum_{s=0}^{k-1}\omega_{w_j}^{s}\wedge
       \omega_{w_\ell}^{k-1-s}\wedge\omega^{n-k}.
\]
The sum includes $s=0$. Factorization and compact Stokes give
\begin{align*}
 &\int_Xh_j\bigl(\omega_{w_j}^{k}-\omega_{w_\ell}^{k}\bigr)
                      \wedge\omega^{n-k}\\
 &\hspace{20mm}=\int_Xh_j\ddc z\wedge T
 =\int_Xz\,\ddc(h_jT).
\end{align*}
Here the full product rule is
\begin{equation}\label{eq:moving-measure-torsion}
 \ddc(h_jT)=\ddc h_j\wedge T+\mathrm{d}h_j\wedge\dc T
             -\dc h_j\wedge \mathrm{d}T+h_j\ddc T.
\end{equation}
The last term is retained. It need not vanish on a Hermitian manifold.

We estimate the four terms. Since $\ddc h_j=\omega_{\alpha_j}-\omega_{a_j}$, the absolute contribution of the first is at most
\[
 \int_X|z|(\omega_{\alpha_j}+\omega_{a_j})\wedge T.
\]
It tends to zero uniformly as $j,\ell\to\infty$ by \cref{lem:bounded-mixed-capacity}, since $w_j-w_\ell\to0$ in capacity. Differentiating each factor shows
$\mathrm{d}T=\mathrm{d}\omega\wedge R$ and $\dc T=\dc\omega\wedge R$, where $R$ is a positive sum of mixed products of degree $(n-2,n-2)$. Moreover $\ddc T$ is a sum of terms containing $\ddc\omega$ and mixed products of degree $n-2$, or $\mathrm{d}\omega\wedge\dc\omega$ and mixed products of degree $n-3$. Terms with negative degree are absent. These smooth torsion factors are bounded in absolute value by fixed multiples of $\omega^2$ and $\omega^3$. As $h_j$ is uniformly bounded, the contribution of $h_j\ddc T$ also tends uniformly to zero by the same capacity estimate.

For the two gradient terms, the pointwise Cauchy--Schwarz inequality for a positive mixed product gives, for example,
\[
 \left|\int_Xz\,\mathrm{d}h_j\wedge\dc\omega\wedge R\right|^2
 \leq C\left(\int_X|z|\,\mathrm{d}h_j\wedge\dc h_j\wedge\omega\wedge R\right)
          \left(\int_X|z|\,\omega^2\wedge R\right).
\]
The second factor tends uniformly to zero. The first is uniformly bounded. For completeness, if $p$ is any of the uniformly bounded potentials and $S$ any positive mixed product of degree $n-1$ occurring here, then
\[
 \int_X\mathrm{d}p\wedge\dc p\wedge S
 =\frac12\int_Xp^2\ddc S-\int_Xp\,\ddc p\wedge S.
\]
Expand $\ddc S$ with all its torsion terms as above, and write $\ddc p=\omega_p-\omega$. Boundedness of $p$ and \cref{lem:bounded-mixed-capacity} bound both integrals uniformly. Finally,
$\mathrm{d}h_j\wedge\dc h_j\leq2\mathrm{d}\alpha_j\wedge\dc\alpha_j+2\mathrm{d}a_j\wedge\dc a_j$
as positive quadratic forms. This proves the asserted bound. The other gradient term has the same estimate. In dimension one $T=1$ and all torsion terms in \eqref{eq:moving-measure-torsion} vanish.

These bounded integration-by-parts and energy identities may first be written for smooth decreasing regularizations in the exact Hermitian background. The local Bedford--Taylor chain rule, its gradient products and quasi-continuity give the displayed identities for bounded potentials. Equivalently they follow from the bounded current calculus of \cref{sec:psh-bounded-products}. Only the common bound on the potentials enters the estimates.

We have shown that the difference of the two $h_j$ integrals tends uniformly to zero when $j,\ell$ are large. For a fixed sufficiently large $\ell$, \cref{lem:fixed-npp-integrals} gives
\[
 \int_Xh_j\,\omega_{w_\ell}^{k}\wedge\omega^{n-k}\longrightarrow0.
\]
Then let $\ell\to\infty$. Together with the estimate for $\alpha_j-a$, this proves the result. If the $a_j$ have a larger fixed lower curvature bound, enlarge $\omega$ to $\Omega$ so that they are $\Omega$-plurisubharmonic. Apply the proved statement with $\Omega$; positivity bounds the original mixed measures by their $\Omega$ counterparts.
\end{proof}

\subsection{Smoothing the measure}
\label{subsec:psh-measure-smoothing}

Coordinate convolution gives smooth approximations of a finite measure. For a varying quasi-plurisubharmonic integrand, the submean inequality transfers the error to the original measure. This is the additional control needed beyond weak convergence of the measures themselves.

For a smooth background $\theta$ with $\PSH(X,\theta)\ne\varnothing$ and a non-pluripolar probability measure $\nu$, put
\[
 A_\theta(\nu)=\sup\left\{\int_X(-w)\,\mathrm{d}\nu:
 w\in\PSH(X,\theta),\ \sup_Xw=0\right\}.
\]
This number may be infinite. Bounded Bedford--Taylor measures satisfy the following estimate.

\begin{lemma}\label{lem:bounded-carrier-integrability}
Let $\omega$ be Hermitian, let $\psi\in\PSH(X,\omega)$ be bounded, and put $\nu=(\omega+\ddc\psi)^n$. There is a finite constant $C$, depending on the fixed backgrounds and $\|\psi\|_\infty$, such that
\[
 \int_X(-w)\,\mathrm{d}\nu\leq C
 \qquad(w\in\PSH(X,\theta),\ \sup_Xw=0).
\]
\end{lemma}
\begin{proof}
On finitely many nested coordinate balls $B'_a\Subset B_a$ covering $X$, choose fixed smooth corrections $q_a,r_a$ with $\ddc q_a\geq\omega$ and $\ddc r_a\geq\theta$, and shift $r_a$ to be nonpositive. Then $\psi+q_a$ is bounded plurisubharmonic, $w+r_a$ is negative plurisubharmonic, and positivity of bounded products gives
\[
 \int_{B'_a}(-w)\,\mathrm{d}\nu
 \leq\int_{B'_a}(-w-r_a)\bigl(\ddc(\psi+q_a)\bigr)^n
 \leq C_a\bigl(1+\|w\|_{L^1(B_a)}\bigr).
\]
The last step is the weighted Chern--Levine--Nirenberg estimate in \cref{sec:psh-local-capacity}. Sum over the cover and use the normalized $L^1$ bound of \cref{thm:normalized-compactness}.
\end{proof}

Choose finitely many coordinate balls $B_a$, smooth functions $\eta_a\geq0$ with compact support in these balls and $\sum_a\eta_a=1$, and inner neighborhoods of their supports whose closures lie in $B_a$. Let $\chi\geq0$ be a smooth radial convolution kernel on the unit ball of $\mathbb C^n$, of integral one. Fix $\epsilon_j\downarrow0$ small enough for every convolution below to stay in the corresponding coordinate ball. Denote convolution of a function in that chart by $S_{a,j}$. For a Borel test function $F$ and a probability measure $\nu$, define a smooth probability measure $\nu_j$ by convolving each coordinate measure $\eta_a\nu$ and adding its extensions by zero. Its defining dual identity is
\begin{equation}\label{eq:measure-smoothing-duality}
 \int_X F\,\mathrm{d}\nu_j
 =\sum_a\int_X\eta_a S_{a,j}F\,\mathrm{d}\nu.
\end{equation}
The kernel is even, so the same convolution occurs on the right. This identity first holds for nonnegative Borel functions, and then for integrable signed functions. In particular $\nu_j\rightharpoonup\nu$.

\Needspace{20\baselineskip}
\begin{lemma}\label{lem:measure-smoothing-integrals}
If $\nu$ is non-pluripolar and $A_\theta(\nu)<\infty$, then
\[
 \sup_j A_\theta(\nu_j)<\infty.
\]
More generally, let $z_j,z$ be quasi-plurisubharmonic functions with a common upper bound and a common smooth lower curvature bound. Suppose they converge both in $L^1(X)$ and in $L^1(\nu)$. Then
\begin{equation}\label{eq:convolution-moving-integrals}
 \int_X|z_j-z|\,\mathrm{d}\nu_j\longrightarrow0.
\end{equation}
Consequently, for every $\lambda>0$,
\[
 \mathrm{e}^{\lambda z_j}\nu_j\rightharpoonup \mathrm{e}^{\lambda z}\nu.
\]
\end{lemma}
\Needspace{9\baselineskip}
\begin{proof}
On each coordinate ball choose a fixed smooth function $q_a$ whose Levi form dominates $\theta$. If $w\in\PSH(X,\theta)$ and $\sup_Xw=0$, then $w+q_a$ is plurisubharmonic. The submean inequality gives
\[
 -S_{a,j}w\leq-w+S_{a,j}q_a-q_a\leq-w+C.
\]
The constant is independent of $w$ and $j$. Integrating \eqref{eq:measure-smoothing-duality} proves $A_\theta(\nu_j)\leq A_\theta(\nu)+C$.

For the second assertion start with any subsequence, and extract so that $z_j\to z$ almost everywhere for smooth volume. The upper envelopes
\[
 H_j=\left(\sup_{k\geq j}z_k\right)^*
\]
decrease to $z$ by \cref{thm:normalized-compactness} and this almost-everywhere convergence. Enlarge the fixed corrections $q_a$ so that $H_j+q_a,z_j+q_a,z+q_a$ are all plurisubharmonic. Since $H_j\geq z_j,z$, the dual identity and the submean inequality imply
\[
 \begin{split}
 \int_X|z_j-z|\,\mathrm{d}\nu_j
 &\leq\sum_a\int_X\eta_a S_{a,j}(2H_j-z_j-z)\,\mathrm{d}\nu\\
 &\leq2\sum_a\int_X\eta_a
       \bigl(S_{a,j}(H_j+q_a)-(z+q_a)\bigr)\,\mathrm{d}\nu
       +\int_X(z-z_j)\,\mathrm{d}\nu.
 \end{split}
\]
All integrals are well defined: each function on the left of a submean inequality is bounded above, and its negative part is controlled by that of its integrable unsmoothed counterpart plus a constant.

Radial averages of a plurisubharmonic function decrease when the radius decreases. Therefore $S_{a,j}(H_j+q_a)$ decreases to $z+q_a$ on the support of $\eta_a$. To identify the limit, fix an index $k$, bound $H_j\leq H_k$ for $j\geq k$, let $j\to\infty$ in the radius, and then let $k\to\infty$. The reverse inequality is the submean inequality. The nonnegative difference in the last display is dominated by an integrable constant minus $z$ on this compact support. Thus its integral tends to zero. The final integral tends to zero by $L^1(\nu)$ convergence. The subsequence criterion proves \eqref{eq:convolution-moving-integrals} for the original sequence.

The exponential is Lipschitz on the common upper range, so its moving $L^1$ error also tends to zero. If $\zeta>0$ is smooth, then $\zeta \mathrm{e}^{\lambda z}=\mathrm{e}^{\lambda z+\log\zeta}$ is bounded and quasi-plurisubharmonic. For any fixed bounded quasi-plurisubharmonic function $F$, the same local submean argument and dominated convergence in \eqref{eq:measure-smoothing-duality} give
\[
 \int_XF\,\mathrm{d}\nu_j\longrightarrow\int_XF\,\mathrm{d}\nu.
\]
Apply this to $F=\zeta \mathrm{e}^{\lambda z}$. Subtracting two positive smooth tests proves the required convergence for every smooth test function. Uniformly bounded masses and uniform approximation of continuous functions by smooth ones prove weak convergence.
\end{proof}

\Needspace{25\baselineskip}
